\section{Definitions and Threat Model}

\begin{table*}[t]
\caption{Core terms.}
\label{tab:core-terms}
\centering
\small
\begin{tabularx}{\textwidth}{L{0.20\textwidth} Y}
\toprule
\textbf{Term} & \textbf{Operational meaning} \\
\midrule
Agent-native repository & A repository structured so a human or agent can discover ownership, legal edit scope, proof lanes, and evidence requirements from versioned local artifacts. \\
Proof receipt & A durable record of a proof command, tool identity, inputs, result, and artifact path. \\
Merge witness & A versioned artifact binding changed paths, owner route, proof receipts, score delta, missing-evidence decision, commit identity, and tool version. \\
Changed-path scope & The concrete path inventory and risk class used to select owners and proof lanes. \\
Owner route & The path-to-owner mapping used to decide who or which policy cell is accountable for review and repair. \\
Generated zone & An output path repaired by changing the declared source contract or generator and regenerating, not by hand edit. \\
Score delta & The difference between baseline and current audit posture; useful for prioritization, but not itself conformance. \\
Repair queue & Ordered findings with rule IDs, owners, proof lanes, rerun commands, evidence requirements, and bounded next actions. \\
Conformance level & A pass/fail claim that required artifacts and gates for HL1--HL5 are current and reproducible. \\
Reference profile & A non-normative stack or implementation profile that can satisfy the standard without defining it. \\
\bottomrule
\end{tabularx}
\end{table*}

\emph{Vibe coding} is any development mode where plausible code is accepted without explicit ownership, proof routing, generated-boundary checks, security gates, and repair evidence. A \emph{continuous audit} is the local, pull-request, and CI loop that reruns those checks at merge time; it is not assumed to be a hosted daemon. The current implementation records an appendable, versioned evidence history of proof plans, command receipts, evidence indexes, audit reports, and score-history rows. This paper reserves stronger ``append-only ledger'' and transparency-log language for future signed receipt work.

The threat model starts with false plausibility. Agents can produce code that looks idiomatic, compiles in isolation, and still violates architecture, security, data ownership, or runtime assumptions. DORA's amplifier finding matters here: AI improves strong systems and worsens weak ones \cite{doraStateAIAssistedSoftwareDevelopment2025}. The repo must therefore treat generated code as a hypothesis until proof lanes accept it.

\begin{table*}[t]
\caption{Threat model for auditable merge.}
\label{tab:threat-model}
\centering
\small
\begin{tabularx}{\textwidth}{L{0.16\textwidth} Y Y}
\toprule
\textbf{Scope} & \textbf{In this paper} & \textbf{Jankurai control} \\
\midrule
Assets & Source code, generated artifacts, CI config, proof receipts, audit history, policy files, owner maps, test maps, and generated-zone manifests. & Versioned manifests, generated-zone declarations, stable rule IDs, and report schemas. \\
Actors & Human maintainer, coding agent, CI runner, malicious contributor, compromised dependency/tool, and stale generated artifact. & Bounded authority, owner routing, source-trust hierarchy, and proof receipts. \\
Trust boundaries & Local workspace, CI, dependency registry, generated-code boundary, repository policy boundary, and release baseline. & Path scope, generated-zone rules, proof lanes, audit reports, and witness decisions. \\
In scope & Missing tests, stale generated files, owner bypass, unproven generated-code edits, unsafe migrations, security-relevant config drift, prompt injection, and false-green tests. & HLT rule families, hard caps, repair queue, security/contract/db/web lanes, and missing-evidence decisions. \\
Out of scope & Proving full semantic correctness, malicious compiler/runtime compromise, and organization-level governance outside the repository. & Explicit limitations, signed-receipt future work, and no claim of total program proof. \\
\bottomrule
\end{tabularx}
\end{table*}

The standard does not claim every repository using AI must adopt Jankurai. It claims a narrower and testable thing: a repository claiming Jankurai conformance must expose enough machine-readable structure for the claim to be audited. Conformance is a public interface, not a vibe. A repo that mentions Jankurai in prose but cannot emit current audit and proof artifacts is not conformant.
